% ==========================================================================
% GENERATED by scripts/split_paper.py. Do not edit this file.
%
% Source: imported_papers/grokking/iclr2027_conference_v2.tex
% Part 8: Circuit measurements and algorithm checks
%
% Edits belong upstream, in the published paper. To change what the
% chapter includes, change imports/grokking.toml and re-run
% the splitter.
% ==========================================================================
\section{Circuit measurements and algorithm checks}
\label{app:measures}

\subsection{Do outgoing norms reflect the distribution of computation?}
\label{app:functional}
We use the participation ratio of outgoing norms rather than a Fourier-based sparsity measure because it needs no Fourier basis, applies unchanged to every task we study, and is invariant to a uniform rescaling of the outgoing weights. It is, however, a structural measure. The network's function is unchanged when a ReLU neuron's incoming weights and bias are multiplied by a positive constant and its outgoing weights are divided by the same constant, so outgoing norms could in principle change without any change in the computation. We check the redistribution with two measures that are tied to the computation.

The first weights each neuron's outgoing norm by its mean absolute activation, $\mathbb{E}_x|h_j(x)| \, \|\mW^{\mathrm{out}}_{:,j}\|_2$, where $h_j(x)$ is the activation of neuron $j$ on input $x$ and the mean runs over the test inputs; this product is invariant to the compensating rescaling above. The second removes each neuron $j$ in turn and records $|\mathbb{E}_x[z_y(x) - z_y^{(-j)}(x)]|$, the absolute value of the mean change of the correct-class logit $z_y$ over the same inputs, where $z_y^{(-j)}$ is the logit with neuron $j$ removed; we apply Equation~\ref{eq:pr} to these ablation effects. Over eight seeds, activation-weighted participation increases from $0.343$ under the reference to $0.813$ under full credit equalization, and ablation participation from $0.422$ to $0.801$. The mean absolute change of the logit, $\mathbb{E}_x|z_y(x) - z_y^{(-j)}(x)|$, and of the margin give the same separation: $0.34$ to $0.49$ under the reference against $0.76$ to $0.98$ under equalization for the logit, and $0.38$ to $0.48$ against $0.85$ to $0.97$ for the margin (three seeds).

Low-norm neurons account for only part of the difference. We recompute participation among the neurons whose outgoing norm is at least one tenth of the layer mean. The reference has $84$ neurons below this threshold and participation $0.889\,(0.006)$ among the rest; full equalization has no neuron below the threshold and participation $0.991\,(0.001)$. Under equalization no neuron falls below the threshold, and the norms above it are also more uniform. It does not make the neurons less selective: the mean fraction of a neuron's input-weight Fourier power at its strongest frequency rises from $0.667$ under the reference to $0.884$ at $\beta = 1$ (five seeds, the experiment of Table~\ref{tab:suff}).

\subsection{Training fit and output scale}
\label{app:norm}
Participation could change because the training fit or the output scale changes. We therefore report the parameter norm $\|\theta\|$ and the logit margin $\gamma = \mathbb{E}_x[z_y(x) - \max_{c \neq y} z_c(x)]$, where $z_c(x)$ is the logit of class $c$ for input $x$, $y$ is the correct label of $x$ and the mean runs over the training set, for the reference, full credit equalization and DFC (Table~\ref{tab:norms}). The reference and the equalized network have nearly the same norm and margin despite their different participation, so the redistribution is not accompanied by a change in fit or output scale. DFC has a much smaller norm and margin: its outputs live on a different scale, which is one reason we use the credit intervention rather than DFC for the matched comparison.

\begin{table}[!ht]
\caption{Parameter norm and logit margin at matched weight decay (modular addition, $p=31$, four seeds; means and standard deviations). All runs have test accuracy $\geq 0.97$. Credit equalization leaves both quantities close to the reference values while producing a circuit at least as distributed as DFC's.}
\label{tab:norms}
\begin{center}
\begin{tabular}{lcc}
\toprule
Rule & $\|\theta\|$ & margin $\gamma$ \\
\midrule
Reference (backpropagation with Adam) & $26.10 \pm 0.15$ & $5.02 \pm 0.02$ \\
Credit equalization ($\beta{=}1$) & $26.65 \pm 0.10$ & $5.15 \pm 0.03$ \\
DFC                & $13.09 \pm 0.02$ & $0.54 \pm 0.00$ \\
\bottomrule
\end{tabular}
\end{center}
\end{table}

\subsection{Checks of the modular-addition computation}
\label{app:algo}
Prior work derives the computation that grokked networks use for modular addition: the first layer represents each operand by a few frequencies $k$, hidden neurons combine them into $\cos$ and $\sin$ of $k(a+b)$, and the logit for candidate $c$ sums $\cos(2\pi k(a+b-c)/p)$ over frequencies, which peaks at $c = a + b$ \citep{nanda2023progress,gromov2023grokking,chughtai2023toy,morwani2024feature}. We check whether the reference, DFC and full credit equalization all implement this computation, and then repeat the checks for the remaining grokking rules.

\paragraph{Fourier measurements.} Let $\vu_j, \vv_j \in \R^p$ be the halves of neuron $j$'s incoming weights that act on the two operands, and let $\hat{u}_j, \hat{v}_j$ be their discrete Fourier transforms. The power of neuron $j$ at frequency $k = 1, \dots, \lfloor p/2 \rfloor$ is
\begin{equation}
s_j(k) \;=\; |\hat{u}_j(k)|^2 + |\hat{v}_j(k)|^2,
\qquad
\hat{u}_j(k) = \textstyle\sum_{a=0}^{p-1} u_j[a]\, e^{-2\pi i k a / p}.
\label{eq:spectrum}
\end{equation}
The aggregate spectrum of a layer sums $s_j(k)$ over neurons.

\paragraph{Frequency ablation.} We keep only the strongest input frequencies, ranked by aggregate power, and count how many are needed to recover $95\%$ of the full test accuracy: $6.7$ for backpropagation with Adam, our example of a concentrated circuit, and $9.3$ for DFC, our example of a distributed circuit. Both networks rely on a small set of input frequencies. A comparison between the reference and full credit equalization, with a random control, follows.

\paragraph{Frequency retention and removal under credit equalization.} We repeat the frequency ablation for the reference and full credit equalization over five seeds and add removal and a random control. Frequencies are ranked once per network by the summed power of its first-layer weights, so the ranking uses no data. We keep or remove the top $k$ frequencies for every $k$ from $0$ to $15$, always keep the constant component, and filter both operand blocks alike. The control keeps or removes $k$ frequencies drawn at random, averaged over $20$ draws per network. For both rules, keeping the top frequencies preserves more test accuracy than keeping random ones, and removing them costs more than removing random ones, at every $k$ from $2$ to $14$ in every seed (Table~\ref{tab:freqcurve}, Figure~\ref{fig:freqcurve}). The equalized network depends on its leading frequencies as the reference does, and spreads that dependence over more of them. Fitted on the training pairs only, the trigonometric function alone classifies every held-out pair correctly in every seed for both rules, with a mean margin between $6.1$ and $6.8$.

\begin{table}[!ht]
\caption{Test accuracy after keeping or removing $k$ input frequencies, ranked by first-layer weight power (top) or drawn at random ($20$ draws per network), modular addition at $p = 31$. Means over five seeds; the complete curves are in Figure~\ref{fig:freqcurve}.}
\label{tab:freqcurve}
\begin{center}
\small
\begin{tabular}{llcccccc}
\toprule
 & & \multicolumn{6}{c}{Frequencies kept or removed, $k$} \\
 & & $2$ & $4$ & $6$ & $8$ & $10$ & $12$ \\
\midrule
Reference & Keep top $k$ & $0.23$ & $0.57$ & $0.86$ & $0.97$ & $1.00$ & $1.00$ \\
 & Keep random $k$ & $0.10$ & $0.21$ & $0.40$ & $0.58$ & $0.76$ & $0.90$ \\
 & Remove top $k$ & $0.83$ & $0.44$ & $0.19$ & $0.06$ & $0.02$ & $0.02$ \\
 & Remove random $k$ & $0.94$ & $0.83$ & $0.62$ & $0.44$ & $0.27$ & $0.13$ \\
\midrule
Equalized & Keep top $k$ & $0.14$ & $0.42$ & $0.69$ & $0.86$ & $0.93$ & $0.96$ \\
 & Keep random $k$ & $0.10$ & $0.23$ & $0.43$ & $0.64$ & $0.80$ & $0.92$ \\
 & Remove top $k$ & $0.91$ & $0.73$ & $0.46$ & $0.23$ & $0.10$ & $0.03$ \\
 & Remove random $k$ & $0.94$ & $0.87$ & $0.71$ & $0.50$ & $0.32$ & $0.16$ \\
\bottomrule
\end{tabular}
\end{center}
\end{table}

\begin{figure}[!ht]
\begin{center}\includegraphics[width=0.8\linewidth]{figures/fig_review_freq_ablation.pdf}\end{center}
\caption{Test accuracy after keeping (left) or removing (right) $k$ input frequencies, ranked by first-layer weight power (solid) or drawn at random (dashed). Means over five seeds, bands one standard deviation across seeds.}
\label{fig:freqcurve}
\end{figure}

\paragraph{Trigonometric fit.} We regress the full logit tensor onto the $(a+b-c)$ trigonometric basis of the Fourier algorithm. The fit explains $0.57$ of the logit variance for the reference, $0.42$ for full credit equalization and $0.90$ for DFC. The residual of the reference agrees across initializations with a correlation of only $0.235$, so there is little evidence of a second computation shared across seeds.

\paragraph{Agreement across runs.} With a fixed data split, seeds of the same rule differ only in their initialization, and their agreement is the yardstick for comparisons across rules. Logit tensors of reference seeds correlate at $0.446$; correlations across rules lie between $0.427$ and $0.467$. Linear centered kernel alignment between the first-hidden-layer activations of the reference and of full equalization is $0.69$, compared with $0.57$ between reference seeds and $0.74$ between equalized seeds. Agreement across rules is thus comparable to agreement across seeds of one rule.

\paragraph{The published alternative.} The pizza algorithm of \citet{zhong2023clock} multiplies the clock logit by an envelope that depends on the difference of the operands, $\lvert\cos(\pi k(a-b)/p)\rvert \cdot \cos(2\pi k(a+b-c)/p)$. Its basis overlaps with the trigonometric one: on a synthetic pizza tensor the trigonometric basis explains $0.81$ of the variance, and on a synthetic clock tensor the pizza basis explains $0.81$, so explained variance alone does not tell the two apart. Two statistics do. Adding the pizza basis to the trigonometric basis raises the explained variance by at most $0.0012$ in every seed of every rule, against $0.19$ for a synthetic pizza tensor and $0.001$ for noise. The distance irrelevance of \citet{zhong2023clock}, the variability of the correct logit at a fixed operand difference relative to its total variability, lies between $0.976$ and $0.990$ for every seed of every rule; they report $0.85$ for a clock network and $0.17$ for a pizza network, and classify networks below $0.6$ as pizza. None of the rules implements this alternative.

\paragraph{The remaining grokking rules.} We repeat the checks for the other learning rules that grok: the eight optimizers of Table~\ref{tab:zoo}, trained as in Appendix~\ref{app:repro}, and the two constructed rules with a fresh or a fading orthogonal component, trained as in Appendix~\ref{app:generalization}, with three seeds each. As above, the seeds share one data split and differ in their initialization, and a reference trained in the same way provides the yardstick (Table~\ref{tab:algoall}). In every rule the trigonometric basis explains between $0.42$ and $0.68$ of the logit variance per seed, the distance irrelevance lies above $0.97$, and the logit correlation and the linear centered kernel alignment of first-hidden-layer activations with the reference across initializations lie at the level of the reference's agreement with itself. Together, the checks show the same periodic computation across all tested grokking rules and distinguish it from the published alternative.

\begin{table}[!ht]
\caption{Algorithm checks for the remaining grokking rules (modular addition, $p=31$, three seeds, one shared data split). Trig: fraction of logit variance explained by the $(a+b-c)$ basis. Distance irrelevance: the statistic of \citet{zhong2023clock}, above $0.6$ for clock networks and below $0.6$ for pizza networks. Correlation and CKA: agreement with the reference trained in the same way, across initializations; the reference rows give its agreement with itself. The optimizers are trained for $20{,}000$ epochs and the constructed rules for $15{,}000$, each with its own reference. The seeds share one data split, so the values differ slightly from Table~\ref{tab:zoo}, whose seeds use different splits.}
\label{tab:algoall}
\begin{center}
\small
\setlength{\tabcolsep}{3pt}
\begin{tabular}{lcccccc}
\toprule
Rule & Test & Circuit PR & Trig & Dist. irrel. & Correlation & CKA \\
\midrule
\multicolumn{7}{l}{\emph{Optimizers, $20{,}000$ epochs}} \\
Reference & $0.999$ & $0.558$ & $0.535$ & $0.983$ & $0.407$ & $0.590$ \\
LARS & $0.975$ & $0.604$ & $0.550$ & $0.986$ & $0.419$ & $0.592$ \\
RMSProp & $0.995$ & $0.496$ & $0.517$ & $0.984$ & $0.456$ & $0.598$ \\
LAMB & $0.998$ & $0.566$ & $0.463$ & $0.984$ & $0.396$ & $0.573$ \\
Adafactor & $0.998$ & $0.489$ & $0.553$ & $0.983$ & $0.432$ & $0.580$ \\
Lion & $0.998$ & $0.479$ & $0.558$ & $0.983$ & $0.461$ & $0.585$ \\
Muon & $0.993$ & $0.494$ & $0.498$ & $0.980$ & $0.407$ & $0.554$ \\
Shampoo & $0.992$ & $0.529$ & $0.525$ & $0.980$ & $0.339$ & $0.518$ \\
Nero & $0.988$ & $0.986$ & $0.455$ & $0.984$ & $0.435$ & $0.659$ \\
\midrule
\multicolumn{7}{l}{\emph{Constructed rules, $15{,}000$ epochs}} \\
Reference & $0.998$ & $0.579$ & $0.532$ & $0.983$ & $0.399$ & $0.594$ \\
Fresh noise & $0.993$ & $0.571$ & $0.531$ & $0.984$ & $0.467$ & $0.621$ \\
Error-gated bias & $0.984$ & $0.688$ & $0.458$ & $0.980$ & $0.415$ & $0.603$ \\
\bottomrule
\end{tabular}
\end{center}
\end{table}
